% =====================================================================
% Chapter 3 -- Theoretical Background
%
% Outline: outlines/ch3_background.md (decisions 2026-10-22).
% Scope: concepts as published, with citations. Own design: Ch. 5; values: Ch. 6.
% Facts: paper_notes/{voxtell,zscl,cpe_clip,lora_hu,liu2026_vlm_cl_survey}.md.
% Notation: thesis_notes.md §5.
% Every paragraph marked CLAUDE-DRAFT is to be rewritten by the author.
% =====================================================================

\chapter{Theoretical Background}
\label{chap:background}

% CLAUDE-DRAFT
This chapter introduces the concepts on which the method of \cref{chap:method} builds.
\Cref{sec:bg-segmentation} describes volumetric segmentation networks, attention and the free-text promptable model VoxTell.
\Cref{sec:bg-cl} states the continual-learning problem in the notation used throughout the thesis.
\Cref{sec:bg-methods} describes the published methods from which our strategies are adapted, and \cref{sec:bg-metrics} the metrics with which we measure adaptation and forgetting.

% ---------------------------------------------------------------------
\section{Free-Text Promptable Volumetric Segmentation}
\label{sec:bg-segmentation}

\subsection{Encoder--Decoder Networks for Volumetric Segmentation}
\label{subsec:bg-unet}

% CLAUDE-DRAFT
We consider the segmentation of one target structure in a 3D image volume.
A volume $\mat{X}$ assigns an intensity to every voxel~$v$ of its voxel set~$\Omega$, and the label mask $\mat{Y}$ assigns $y_v = 1$ to the voxels of the structure and $y_v = 0$ to all others.
A segmentation network predicts one logit $z_v$ per voxel, the logistic function~$\sigma$ turns it into a foreground probability $\hat{y}_v = \sigma(z_v)$, and a threshold gives the binary prediction~$\hat{\mat{Y}}$.
In a free-text promptable model, the target is not fixed by the output channels of the network but named by a text prompt~$p$, so one network with parameters~$\theta$ computes the logits $\mat{Z} = f_\theta(\mat{X}, p)$ for any pair of volume and prompt.
We write $\mathcal{L}_{\mathrm{seg}}$ for the segmentation loss that compares $\mat{Z}$ with $\mat{Y}$.
VoxTell uses the sum of the binary cross-entropy and the soft Dice loss~\cite[Sec.~3]{rokuss2025voxtellfreetextpromptableuniversal}; the soft Dice term is computed over the whole patch and therefore does not let the many background voxels dominate the loss as a voxel-wise average does~\cite{milletari2016vnet}.
The form used in this thesis is given in \cref{eq:method-seg-loss}.

% CLAUDE-DRAFT
VoxTell builds on the U-Net design in its 3D form~\cite{3d_unet}, which its authors regard as still the strongest architecture for 3D medical image segmentation~\cite[p.~4]{rokuss2025voxtellfreetextpromptableuniversal}.
An encoder computes features at successively lower resolutions in a sequence of stages, and a decoder upsamples them back to the input resolution.
Skip connections pass the encoder features of every stage to the decoder stage of the same resolution, so the decoder can use spatial detail that was lost by downsampling.
We write $\mat{F}_s$ for the encoder features of stage~$s$ and $C_s$ for their number of channels, with $s = 0$ at the full resolution and larger~$s$ at coarser resolutions.
Because a whole volume and its intermediate features do not fit into the memory of a GPU, such networks are trained on patches cut from the volumes and applied at inference with a sliding window, whose overlapping patch predictions are combined into one prediction per volume.
VoxTell is built on the nnU-Net framework~\cite[S-p.~1]{rokuss2025voxtellfreetextpromptableuniversal}, which configures these steps automatically for a given dataset~\cite{isensee2018nnunetselfadaptingframeworkunetbased}, and uses its residual-encoder configuration ResEnc~L~\cite[Sec.~3.1]{isensee2024nnunetrevisited}.

\subsection{Attention}
\label{subsec:bg-attention}

% CLAUDE-DRAFT
Attention lets a set of query vectors collect information from a set of source vectors, weighted by how well each query matches each source~\cite{NIPS2017_3f5ee243}.
Let the rows of $\mat{U} \in \mathbb{R}^{n \times d}$ be $n$ query vectors and the rows of $\mat{S} \in \mathbb{R}^{m \times d}$ be $m$ source vectors of dimension~$d$.
Multi-head attention projects them with learned matrices $\mat{W}_q, \mat{W}_k, \mat{W}_v \in \mathbb{R}^{d \times d}$ to queries, keys and values and splits each projection into $n_h$ heads of dimension $d_h = d / n_h$.
With $\mat{W}_q^{(i)}, \mat{W}_k^{(i)}, \mat{W}_v^{(i)} \in \mathbb{R}^{d_h \times d}$ the rows of the projections that belong to head~$i$,
\begin{equation}
\begin{aligned}
\operatorname{head}_i &= \operatorname{softmax}\!\left( \frac{\mat{U}\mat{W}_q^{(i)\top} \big(\mat{S}\mat{W}_k^{(i)\top}\big)^{\!\top}}{\sqrt{d_h}} \right) \mat{S}\mat{W}_v^{(i)\top},\\
\operatorname{MHA}(\mat{U}, \mat{S}) &= \big[\operatorname{head}_1, \dots, \operatorname{head}_{n_h}\big]\, \mat{W}_o^{\top},
\end{aligned}
\label{eq:bg-attention}
\end{equation}
where the softmax runs over the $m$ sources for each query, $[\cdot]$ concatenates the heads along the feature dimension and $\mat{W}_o \in \mathbb{R}^{d \times d}$ is the output projection.
We write all projections as maps $\vec{x} \mapsto \mat{W}\vec{x}$ of column vectors, which for row-wise inputs gives the transposes in \cref{eq:bg-attention}.
Each query thus receives a weighted mean of the values, with weights given by the similarity of its query to every key.
Attention itself does not depend on the order or position of the sources, so a positional encoding is usually added to them~\cite{NIPS2017_3f5ee243}.

% CLAUDE-DRAFT
In self-attention, the queries and the sources are the same vectors ($\mat{S} = \mat{U}$); in cross-attention, they come from different inputs, for example a text embedding and image features.
Keys and values are computed from each source vector separately, so further rows appended to~$\mat{S}$ give every query additional positions to attend to without any change to the projections, the mechanism that the prompt-tuning strategy of \cref{subsec:method-prompt} relies on.
A transformer decoder layer combines attention blocks with a feed-forward network of two linear layers applied to every query separately, and wraps each block in a residual connection and a layer normalisation~\cite{NIPS2017_3f5ee243}.
In the original design, each decoder layer applies self-attention to its queries and then cross-attention to the encoder output; VoxTell's prompt decoder uses the same building blocks, with the text embedding as query and image features as sources~\cite[Eq.~3]{rokuss2025voxtellfreetextpromptableuniversal}.

\subsection{VoxTell}
\label{subsec:bg-voxtell}

% CLAUDE-DRAFT
VoxTell is a vision--language model that maps a 3D volume and a free-text prompt, from a single word to a clinical sentence, to a segmentation mask of the described structure~\cite[p.~1]{rokuss2025voxtellfreetextpromptableuniversal}.
It was trained on more than 62{,}000 \gls{ct}, \gls{mri} and PET volumes from 158 public datasets, with a vocabulary of 1{,}087 anatomical and pathological concepts~\cite[p.~4]{rokuss2025voxtellfreetextpromptableuniversal}.
The model has four parts (\cref{fig:ch3_voxtell}): a frozen text encoder, a convolutional image encoder, a transformer prompt decoder that relates the text to the image, and a convolutional decoder into which the text is fused at every resolution~\cite[Sec.~3]{rokuss2025voxtellfreetextpromptableuniversal}.

% CLAUDE-DRAFT
The text encoder is the embedding model Qwen3-Embedding-4B~\cite{zhang2025qwen3embedding}, which maps the text prompt~$p$, preceded by a fixed instruction to retrieve the anatomical entity and location it describes, to an embedding $\vec{q} = f_{\mathrm{text}}(p)$~\cite[Eq.~2, S-p.~1]{rokuss2025voxtellfreetextpromptableuniversal}.
The image encoder is the residual encoder of nnU-Net with six stages, $s = 0, \dots, 5$, of 32, 64, 128, 256, 320 and 320 channels~\cite[S-p.~1]{rokuss2025voxtellfreetextpromptableuniversal}; for a patch of $192^3$ voxels, the coarsest stage has $6^3$ voxels.

% CLAUDE-DRAFT
The prompt decoder consists of $L_{\mathrm{dec}} = 6$ transformer decoder layers with eight heads and dimension $d = 2048$~\cite[Sec.~3]{rokuss2025voxtellfreetextpromptableuniversal}.
Its query is the text embedding~$\vec{q}$, and its keys and values are image features: according to the paper the features of the bottleneck~\cite[Eq.~3]{rokuss2025voxtellfreetextpromptableuniversal}, in the released implementation used in this thesis those of stage~4 (\cref{sec:method-adaptation}).
One two-layer perceptron per decoder stage maps the output of the prompt decoder to a text-guidance matrix $\mat{E}_s \in \mathbb{R}^{G \times C_s}$, that is, $G = 32$ guidance vectors with the channel number of that stage~\cite[Sec.~3, S-p.~1]{rokuss2025voxtellfreetextpromptableuniversal}.
Because the query has attended to the image before it is mapped to the guidance, the guidance depends on both the prompt and the image.

% CLAUDE-DRAFT
The decoder upsamples the features stage by stage and fuses the guidance into every stage~\cite[Eqs.~4--6]{rokuss2025voxtellfreetextpromptableuniversal}.
Starting from $\mat{H}_5 = \mat{F}_5$, it computes for $s = 4, \dots, 0$
\begin{equation}
\mat{H}'_s = \operatorname{ConvBlock}\big(\big[\operatorname{Up}(\mat{H}_{s+1});\, \mat{F}_s\big]\big),
\qquad
\mat{H}_s = \big[\mat{H}'_s;\, \mat{E}_s \odot \mat{H}'_s\big],
\label{eq:bg-voxtell-fusion}
\end{equation}
where $\operatorname{Up}$ upsamples to the resolution of stage~$s$, $[\cdot\,;\cdot]$ concatenates along the channel dimension, and $\mat{E}_s \odot \mat{H}'_s$ denotes the $G$ maps of dot products between each guidance vector and the $C_s$-dimensional feature vector at every voxel.
Each guidance vector thus marks the voxels whose features match it, and these $G$ maps are available as additional channels to all following convolutions.
The authors contrast this multi-stage fusion with the single late fusion of earlier text-promptable models such as SAT~\cite{zhao2025sat}, in which the text meets the image features only once, at the end of the decoder.
They argue that late fusion forces the image features to be computed without knowledge of the prompt and that prompts naming a location need repeated interaction between text and image~\cite[p.~4]{rokuss2025voxtellfreetextpromptableuniversal}.

% CLAUDE-DRAFT
VoxTell was trained with the sum of the soft Dice and the binary cross-entropy loss on the outputs of all five decoder stages, weighted from 1 at the full resolution down to $1/16$ as in nnU-Net (deep supervision)~\cite[Eq.~7, S-p.~1]{rokuss2025voxtellfreetextpromptableuniversal}.
Training used patches of $192^3$ voxels, 85\,\% of them centred on foreground, the data augmentation of nnU-Net without left--right mirroring, and SGD with an initial learning rate of $10^{-4}$ and polynomial decay~\cite[S-p.~1]{rokuss2025voxtellfreetextpromptableuniversal}.
Each training image was paired with three prompts, two naming structures present in the image and one naming an absent structure, so that the model also learns to return an empty mask; a prompt was the default name of its concept with probability 0.25 and a rephrased variant otherwise~\cite[S-p.~1]{rokuss2025voxtellfreetextpromptableuniversal}.
\Cref{fig:ch3_voxtell} shows the model as it is used in this thesis; \cref{sec:method-adaptation} describes where the released implementation differs from the paper.

\begin{figure}[tb]
\centering
\includegraphics{figures/ch3_voxtell.pdf}
% CLAUDE-DRAFT
\caption[Architecture of VoxTell]{Architecture of VoxTell (adapted from~\cite[Fig.~2]{rokuss2025voxtellfreetextpromptableuniversal}).
The image encoder computes features at six stages and passes them to the decoder stage of the same resolution.
The frozen text encoder (lock) embeds the text prompt into the query~$\vec{q}$, which the prompt decoder relates to the image features; its output is mapped to the guidance~$\mat{E}_s$, which is fused into each of the five decoder stages (circles, \cref{eq:bg-voxtell-fusion}).
Drawn as used in this thesis: the prompt decoder attends to the stage-4 features and only the full-resolution output enters the loss, whereas the paper describes attention to the bottleneck and deep supervision.
Input: a patch of SKM-TEA volume MTR\_005 (test split) with the text prompt ``femoral cartilage''; the output panel shows the reference mask for illustration.}
\label{fig:ch3_voxtell}
\end{figure}

% ---------------------------------------------------------------------
\section{Continual Learning}
\label{sec:bg-cl}
% Scope (decided 2026-10-22): only the problem in symbols and the relation of our setting to the
% scenarios. Informal concepts: Ch. 1; scenario definitions and method families: Ch. 2.

% CLAUDE-DRAFT
In continual learning, a model is trained on a sequence of tasks $t = 1, \dots, T$, one after the other.
Training on task~$t$ starts from the parameters $\theta_{t-1}$ reached after the previous task and uses the training data $\mathcal{D}_t$ of task~$t$, which yields the parameters~$\theta_t$; the data of earlier tasks are not available again unless a strategy keeps part of them.
In our setting, the sequence starts from the pretrained parameters~$\theta_0$, and each task comes with its own text prompt~$p_t$.
We call the training on one task a stage, so stage $j = 0$ denotes the pretrained model and stage $j = t$ the model after training on task~$t$.

% CLAUDE-DRAFT
To see how training on later tasks affects earlier ones, the model is evaluated on every task after every stage~\cite[Fig.~4b]{zheng2023preventingzeroshottransferdegradation}.
We write $a_{i,j}$ for the performance on task~$i$ after stage~$j$, here the mean \gls{dsc} on the test set of task~$i$, with $i = 1, \dots, T$ and $j = 0, \dots, T$; together, these values form the evaluation matrix.
Its diagonal entries $a_{i,i}$ measure how well a task is learned when it is trained, the entries $a_{i,j}$ with $j > i$ how much of it is kept afterwards, and the entries with $j < i$ how the model performs on a task before training on it.
Catastrophic forgetting is the case in which training on later tasks lowers the performance on earlier ones~\cite{french1999catastrophic}, that is, $a_{i,j} < a_{i,i}$ for $j > i$.
A strategy that changes the parameters little keeps the earlier entries but may learn new tasks poorly, and a strategy that adapts freely does the opposite; this conflict between stability, the retention of previous knowledge, and plasticity, the integration of new knowledge, is known as the stability--plasticity dilemma~\cite[Sec.~1]{delange2022}.
For a pretrained model, a further question is whether it keeps the abilities it had before the first task: ZSCL showed that continual training of CLIP degrades its zero-shot transfer to data it has not been trained on, in addition to forgetting the trained tasks~\cite[Sec.~1]{zheng2023preventingzeroshottransferdegradation}.
The metrics that summarise the evaluation matrix are defined in \cref{sec:bg-metrics}.

% CLAUDE-DRAFT
Continual-learning settings are usually classified into task-, domain- and class-incremental scenarios, which differ in what the model is told at inference and what it must infer itself (\cref{sec:cl-vlm}).
Our setting does not fit any of them cleanly.
All tasks share the same images, and each task adds a new target structure with its own masks.
At inference, the text prompt names the structure to segment in every strategy, so the model never has to infer which of the learned targets is wanted, the difficulty that defines class-incremental learning.
At the same time, all strategies except task-specific LoRA use one parameter set for all tasks and receive no task identity, unlike in task-incremental learning.
We therefore describe the strategies by their mechanism and distinguish only whether they use the task identity at inference (task-aware inference, \cref{sec:method-setting}) or one parameter set for all tasks; the scenario names are used only to describe related work in its own terms.

% ---------------------------------------------------------------------
\section{Methods Adapted by the Strategies}
\label{sec:bg-methods}

% CLAUDE-DRAFT
This section describes the published mechanisms on which our strategies build, each in the form in which it was proposed.
How we adapted them to VoxTell, and where our versions differ, is described in \cref{sec:method-strategies}.

\subsection{Low-Rank Adaptation}
\label{subsec:bg-lora}

% CLAUDE-DRAFT
Low-rank adaptation (LoRA) keeps a pretrained weight matrix $\mat{W}_0 \in \mathbb{R}^{d_{\mathrm{out}} \times d_{\mathrm{in}}}$ frozen and learns an additive update of low rank~\cite[Sec.~4.1]{hu2022lora},
\begin{equation}
\mat{W} = \mat{W}_0 + \gamma_r \, \mat{B}\mat{A},
\qquad \mat{B} \in \mathbb{R}^{d_{\mathrm{out}} \times r}, \quad \mat{A} \in \mathbb{R}^{r \times d_{\mathrm{in}}}, \quad r \ll \min(d_{\mathrm{in}}, d_{\mathrm{out}}),
\label{eq:bg-lora}
\end{equation}
where only $\mat{A}$ and $\mat{B}$ are trained and $\gamma_r$ is a scaling factor.
$\mat{A}$ is initialised randomly and $\mat{B}$ with zeros, so the update is zero at the start and training begins from the pretrained model.
After training, $\gamma_r\mat{B}\mat{A}$ can be added to $\mat{W}_0$, so the adapted model has no additional inference cost, and switching to another task only requires exchanging the update~\cite[Sec.~4.1]{hu2022lora}.
Hu et al.\ applied LoRA to the projection matrices of self-attention (\cref{eq:bg-attention}) and, comparing choices under a fixed number of trainable parameters, found that adapting the query and value projections together gave the best performance overall~\cite[Sec.~4.2, 7.1, Table~5]{hu2022lora}.

% CLAUDE-DRAFT
In the original formulation, $\gamma_r = \alpha / r$ with a constant hyperparameter~$\alpha$~\cite[Sec.~4.1]{hu2022lora}.
Kalajdzievski showed that this factor makes the gradients of the adapter collapse as the rank grows, so that large ranks learn no better than small ones, and proved that only a factor proportional to $1/\sqrt{r}$ keeps the outputs and gradients of the adapter stable for all ranks~\cite[Sec.~2.1, 3]{kalajdzievski2023rankstabilizationscalingfactor}.
The resulting rank-stabilised LoRA (rsLoRA) sets
\begin{equation}
\gamma_r = \frac{\alpha}{\sqrt{r}}.
\label{eq:bg-rslora}
\end{equation}
In fine-tuning experiments with a large language model and ranks up to 2048, rsLoRA improved with increasing rank while standard LoRA did not, and the gradient norms showed the same pattern with SGD as with AdamW~\cite[Sec.~4]{kalajdzievski2023rankstabilizationscalingfactor}.

\subsection{Replay}
\label{subsec:bg-replay}

% CLAUDE-DRAFT
Replay methods store samples of earlier tasks, or generate pseudo-samples, and train on them again while learning a new task~\cite[Sec.~3.1]{delange2022}.
In the simplest form, rehearsal, a limited subset of stored training samples is added to the training data of each new task.
Methods differ in which samples they keep and how they divide a fixed memory among the tasks; iCaRL, for example, keeps per class the exemplars that best approximate the class mean in feature space~\cite[Sec.~3.1]{delange2022}.
Because rehearsal trains on a subset of the data of all tasks seen so far, its performance is bounded by joint training on all tasks, and it can overfit the few stored samples~\cite[Sec.~3.1]{delange2022}.
Storing data also conflicts with storage limits and privacy constraints~\cite[Sec.~1]{delange2022}, which in medical imaging restrict what may be kept from earlier patients.

\subsection{Knowledge Distillation and ZSCL}
\label{subsec:bg-zscl}

% CLAUDE-DRAFT
In knowledge distillation, a student model is trained to reproduce the output distribution of a teacher model, softened by a temperature in the softmax~\cite[Sec.~2]{hinton2015distilling}.
Continual learning uses distillation as a regulariser: Learning without Forgetting (LwF) trains on the data of the new task only and uses the outputs of the model from before this task as soft labels for the earlier tasks~\cite{li2017learningforgetting}~\cite[Sec.~3.2.1]{delange2022}.
The teacher is thus the previous model, and the distillation data are those of the current task.

% CLAUDE-DRAFT
ZSCL addresses continual learning of CLIP, for which the problem is not only forgetting of the trained tasks but also the loss of zero-shot transfer to data that it has not been trained on~\cite[Sec.~1]{zheng2023preventingzeroshottransferdegradation}.
It changes two choices of LwF.
The teacher is the initial pretrained model instead of the previous one, and the distillation data are a separate reference dataset of unlabelled public images and arbitrary texts instead of the current task's data~\cite[Sec.~3.2]{zheng2023preventingzeroshottransferdegradation}.
For each image, CLIP's similarities to a set of texts are turned into a distribution by a softmax, and the loss compares these distributions of student and teacher by a cross-entropy; the same is done in the other direction, for each text over the images of a batch, and both terms are added to the task loss with a weight~$\lambda$~\cite[Sec.~3.2, Eqs.~2--4]{zheng2023preventingzeroshottransferdegradation}.
In addition, ZSCL limits the change of the parameters themselves.
Every $I_{\mathrm{we}}$ iterations it updates a weight ensemble, the running mean of the parameters along the training trajectory, starting from $\theta_0$,
\begin{equation}
\hat{\theta}_n = \frac{1}{n+1}\,\theta + \frac{n}{n+1}\,\hat{\theta}_{n-1}, \qquad \hat{\theta}_0 = \theta_0,
\label{eq:bg-we}
\end{equation}
where $\theta$ are the current parameters and $n$ counts the ensemble updates~\cite[Sec.~3.3, Eq.~7]{zheng2023preventingzeroshottransferdegradation}.
The full method further adds a penalty on the squared distance between the parameters and a reference set of parameters~\cite[Sec.~3.3, Eq.~5, Table~2]{zheng2023preventingzeroshottransferdegradation}.

\subsection{Prompt Learning and CPE-CLIP}
\label{subsec:bg-prompt}

% CLAUDE-DRAFT
Prompt learning adapts a frozen vision--language model by training a few input vectors instead of its weights.
CoOp replaces the hand-written context words of a text prompt, such as ``a photo of a'', by learnable context vectors in front of the class name, while all parameters of CLIP stay fixed~\cite{zhou2022coop}~\cite[Sec.~3.2]{zhou2022cocoop}.
CoCoOp adds a small network, the meta-net, that computes a token from the image features and adds it to every context vector, so the prompt depends on the input image; its authors report that this generalises better to classes not seen in training~\cite[Sec.~3.3, 4.1]{zhou2022cocoop}.

% CLAUDE-DRAFT
CPE-CLIP applies prompt learning to few-shot class-incremental learning, in which a base session with many classes and examples is followed by sessions that each add a few classes with few examples, and no task identity is given at evaluation~\cite[Sec.~1, 3]{D_Alessandro_2023}.
It keeps CLIP frozen and learns prompt tokens in the layers of the language encoder up to a given depth, replacing the prompt tokens of each layer by new learned ones; the prompts of the vision encoder are derived from the language prompts by a learned linear projection~\cite[Sec.~4, Eqs.~2--8]{D_Alessandro_2023}.
To limit forgetting, it scales the gradients of the prompts and of the projection in session~$t$ by
\begin{equation}
\beta_t = \frac{|C_t|}{\sum_{\tau=0}^{t} |C_\tau|},
\label{eq:bg-cpe-scaling}
\end{equation}
where $C_t$ is the set of classes of session~$t$, so that the prompts change less as more classes have been learned~\cite[Sec.~4, Eqs.~11--13]{D_Alessandro_2023}.
In the authors' ablation, removing this scaling degraded performance more than removing any other component~\cite[Sec.~5.5, Fig.~2]{D_Alessandro_2023}.

% ---------------------------------------------------------------------
\section{Evaluation Metrics}
\label{sec:bg-metrics}
% Definitions as implemented in scripts/compute_cl_metrics.py (reference implementation).
% Which metrics are reported and why: Ch. 6 setup (decision 2026-10-21).

% CLAUDE-DRAFT
The overlap between a binary prediction $\hat{\mat{Y}}$ and the label mask $\mat{Y}$ is measured with the \gls{dsc}~\cite{dice1945measures},
\begin{equation}
\operatorname{DSC}(\hat{\mat{Y}}, \mat{Y}) = \frac{2\,|\hat{\mat{Y}} \cap \mat{Y}|}{|\hat{\mat{Y}}| + |\mat{Y}|},
\label{eq:bg-dsc}
\end{equation}
where $|\cdot|$ counts the foreground voxels of a mask and $\hat{\mat{Y}} \cap \mat{Y}$ denotes the voxels that are foreground in both.
The \gls{dsc} is 0 for disjoint masks and 1 for identical ones.
Because it is normalised by the sizes of the two masks and not by the size of the volume, the many background voxels do not inflate it, which matters for structures as small as the knee cartilages and menisci (\cref{tab:ch4_structures}).
Each entry $a_{i,j}$ of the evaluation matrix (\cref{sec:bg-cl}) is the mean \gls{dsc} over the test volumes of task~$i$.
The continual-learning metrics below were defined for classification accuracy, but Chaudhry et al.\ note that a task-specific measure such as the intersection over union of a segmentation can take the place of accuracy without changing the definitions~\cite[Sec.~3]{Chaudhry_2018}.

\begin{figure}[tb]
\centering
\includegraphics{figures/ch3_metrics.pdf}
% CLAUDE-DRAFT
\caption[Entries of the evaluation matrix used by each metric]{Entries of the evaluation matrix used by each metric (\cref{eq:bg-adaptation,eq:bg-forgetting,eq:bg-control-drop,eq:bg-final,eq:bg-transfer}).
Rows are the tasks $i = 1, \dots, T$, columns the stages $j = 0, \dots, T$, with stage~0 the pretrained model; the separate row~$c$ is the control task.
Dark cells enter the mean of a metric; light cells are the reference from which the dark cell of the same row is subtracted: the maximum over the light cells for Final Forgetting and the pretrained value $c_0$ for the control-task drop.}
\label{fig:ch3_metrics}
\end{figure}

% CLAUDE-DRAFT
We summarise the evaluation matrix with four metrics from the literature and add a fifth for the control task, a task that is not part of the sequence and whose mean \gls{dsc} after stage~$j$ we write as~$c_j$.
\Cref{fig:ch3_metrics} shows which entries each metric uses.
{\allowdisplaybreaks % local: lets this block break across pages
\begin{align}
\text{Adaptation} &= \frac{1}{T} \sum_{i=1}^{T} a_{i,i},
\label{eq:bg-adaptation}\\
\text{Final Forgetting} &= \frac{1}{T-1} \sum_{i=1}^{T-1} \Big( \max_{j \in \{1, \dots, T-1\}} a_{i,j} - a_{i,T} \Big),
\label{eq:bg-forgetting}\\
\text{Control-task drop} &= c_0 - c_T,
\label{eq:bg-control-drop}\\
\text{Final Performance} &= \frac{1}{T} \sum_{i=1}^{T} a_{i,T},
\label{eq:bg-final}\\
\text{Transfer} &= \frac{1}{T-1} \sum_{i=2}^{T} \frac{1}{i-1} \sum_{j=1}^{i-1} a_{i,j}.
\label{eq:bg-transfer}
\end{align}}
Lower values are better for Final Forgetting and the control-task drop, higher values for the other three metrics.

% CLAUDE-DRAFT
Adaptation measures how well each task is learned when it is trained.
It is the term $a_{k,k}$ of the intransigence measure of Chaudhry et al., which subtracts it from the performance of a reference model trained on the data of all tasks at once~\cite[Sec.~3, Eq.~4]{Chaudhry_2018}; Adaptation reports this term without the reference.

% CLAUDE-DRAFT
Final Forgetting is the average forgetting of Chaudhry et al.\ after the last task: for each earlier task, the best performance at any previous stage minus the performance at the end~\cite[Sec.~3, Eq.~3]{Chaudhry_2018}.
Chaudhry et al.\ index the matrix by stage first; \cref{eq:bg-forgetting} is written in our convention.
Taking the maximum over all previous stages, instead of the value right after training, also counts a loss of gains that a task made while later tasks were trained~\cite[Sec.~3]{Chaudhry_2018}.
As in their definition, the maximum runs over all stages $1, \dots, T-1$; in our matrix this includes stages before task~$i$ was trained, because every task is evaluated at every stage.
A negative value means that training on later tasks improved earlier ones, which Chaudhry et al.\ call positive backward transfer~\cite[Sec.~3]{Chaudhry_2018}.

% CLAUDE-DRAFT
The control-task drop follows the control set of SPU, which measures how much of a pretrained model's general ability is lost during continual learning as the accuracy of the frozen pretrained model on the control set minus that of the continually trained model~\cite[Sec.~5.1, Eq.~6]{Zhang_2024_CVPR}.
We compute it with the \gls{dsc} for one task sequence, between the pretrained model and the model after the last stage, so a positive value is a loss.
SPU samples its control set from the pretraining distribution~\cite[Sec.~3]{Zhang_2024_CVPR}; our control task instead comes from a dataset that VoxTell was not trained on, with a concept that is well represented in its training data (\cref{sec:data-aeropath}).

% CLAUDE-DRAFT
Final Performance is the metric Last of ZSCL, the mean performance on all tasks after the last stage~\cite[Sec.~4, Fig.~4b]{zheng2023preventingzeroshottransferdegradation}; it describes the model that a continual-learning run delivers.
Transfer is the metric of the same name in ZSCL: the mean performance on each task at the stages before it is trained~\cite[Sec.~4, Fig.~4b]{zheng2023preventingzeroshottransferdegradation}.
For a pretrained model, it combines the ability that the model brings to a structure before seeing its labels with the effect that training on other structures has on it.

